\chapter*{Glossary}
\addcontentsline{toc}{chapter}{Glossary}

\begin{description}\setlength{\itemsep}{4pt}
\item[Aspect] A retail-relevant topic drawn from a fixed eight-item taxonomy
      (pricing, product quality, customer service, store experience,
      online/app, delivery/pickup, returns, account/login) used to tag
      opinions in each post.
\item[Curriculum learning] A staged training scheme in which a model is first
      exposed to easy or general examples and then progressively harder or
      more domain-specific ones.
\item[Kanban lifecycle] A four-state workflow board (Triaged --- Acknowledged
      --- In Progress --- Resolved) used by analysts to track P1/P2 posts.
\item[Multi-pass captioning] A vision pipeline that runs the caption model
      twice --- first for a free-form description, then for constrained OCR
      and retail-signal extraction --- to reduce hallucination and improve
      text recovery.
\item[Priority tier] A rule-derived severity label (P1: highest, P2:
      secondary) applied to trust-weighted negative posts to focus analyst
      attention.
\item[Trust score] A 0--1 score combining account-metadata heuristics and an
      LLM-graded credibility signal; used with model confidence as an
      admission gate before aggregation.
\item[Zero-shot NLI] Using a Natural Language Inference model to score
      whether a hypothesis (``This post is about pricing.'') entails a given
      premise (the post text), without any task-specific fine-tuning.
\end{description}
\cleardoublepage
